\chapter{Anillos y grupo (multiplicativo y cíclico)}
\label{anexo:iso}

\section{Grupo y Anillo}
Se llama grupo al par (G,+) donde \\
G es un conjunto \\
+ : G X G → G \quad es una ley de composición interna sobre G, asociativa, con neutro e inverso respecto a +. \\
El orden de un grupo es el número de elementos de G. \\
Para cada entero $p$ hay un grupo denominado el grupo cíclico de orden p, donde: \\
G = $Z_p$ = \{ 0, 1, …, (p-1) \} y la operación + es la suma módulo $p$: $i + j = (i + j) \mod p$. \\
Se llama grupo abeliano a un grupo conmutativo, es decir, a un grupo donde la operación + es conmutativa. \\

Se llama anillo (ring) a una tripleta (R,+,×) donde \\
(R,+) \quad es un grupo abeliano \\
× : R X R → R \quad es una ley de composición interna sobre R, asociativa \\
× es distributiva sobre +:	a × (b + c) = (a × b) + (a × c) \\
Es usual denominar adición y multiplicación a las operaciones + y ×; y a sus neutros 0 y 1. \\

Si la operación × es conmutativa y tiene neutro se agrega al nombre de anillo: conmutativo, con unidad. \\
Si además cada elemento distinto del 0 tiene inverso multiplicativo se agrega al nombre de anillo: con división (o cuerpo o campo o field).

\section{Clases de congruencia módulo $n$ y algo de notación}
Recordemos que la relación de congruencia define una relación de equivalencia, de modo que el conjunto de los enteros $\Z$ queda particionado en clases de equivalencia. \\
La clase de equivalencia de a es: \\
$[a]_n = \{ b \in \Z \; / \; a \equiv b \pmod n\} = \{\dots, a - 2n, a - n, a, a + n, a + 2n, \dots \}$ \\
El conjunto de todas las clases se denota por $\Z/n\Z$ y es $\Z/n\Z = \{[0]_n, [1]_n, ..., [n -1]_n \}$. \\
 Podemos hacer esta interpretación: tomamos el conjunto de los enteros $\Z$, consideramos los múltiplos de $n$, es decir, $n\Z$ y $\Z/n\Z$ (que suele recibir el nombre de grupo cociente) son los enteros agrupados por su resto módulo $n$. \\
 Es decir, agrupamos los elementos que difieren en un múltiplo de $n$ en la misma clase de equivalencia. Así, $\Z/n\Z$ es el conjunto de clases de enteros que tienen el mismo resto al dividir por $n$, de modo que dos enteros se identifican si difieren en un múltiplo de $n$. \\
Es obvio que análoga interpretación reciben $\Z/p\Z$, $\Z/q\Z$, etc. 

\section{Las operaciones}
Dadas las clases $[a]_n$ y $[b]_n$, que son conjuntos, se define las operaciones suma y producto entre estos conjuntos así: \\
$[a]_n + [b]_n : = [a + b]_n$ \\
$[a]_n \cdot [b]_n : = [a \cdot b]_n$ \\
A diferencia de una operación entre enteros: \\
$(a+b) \mod n = ( (a \mod n) + (b \mod n) ) \mod n$ \\
$(a \cdot b) \mod n = ( (a \mod n) \cdot (b \mod n) ) \mod n$ \\
nótese que la operación módulo está grabada en la identidad misma de la clase al incluirla, por definición, en la congruencia.  \\
"Reiteramos que las clases son conjuntos. Aún así, es usual trabajar directamente con sus representantes canónicos: el elemento $0$ para la clase $[0]_n$, y así sucesivamente hasta $n-1$ para $[n-1]_n$, identificando por comodidad al conjunto a través de su residuo mínimo no negativo.

\begin{ejemplo}
Con $a=9, b=11, n=5$. \\
En los enteros: \\
$(a+b) \mod n = (9+11) \mod 5 = 20 \mod 5 = ( (9 \mod 5) + (11 \mod 5) ) \mod 5 = ( (4) + (1) ) \mod 5 = 5 \mod 5 = 0$. \\
En $\Z/n\Z$: \\
$[a]_n + [b]_n = [9]_5 + [11]_5 = [9 + 11]_5 = [20]_5$ \\
$= [4]_5 + [1]_5 = [0]_5$. \\
Note que en la última línea, por comodidad, hemos anotado los representantes canónicos de las clases de equivalencia del 9, 11 y 20.
\end{ejemplo}

Es usual mostrar que las operaciones están bien definidas. Es decir, que las mismas no dependen de los elementos elegidos.\\
Suma: \\
Supongamos que tomamos otros representantes $a' \in [a]_n$ y $b' \in [b]_n$. \\  
Eso significa: \\
$a' \equiv a \pmod{n}, \quad b' \equiv b \pmod{n}$. \\
Entonces:
$a' = a + kn, \quad b' = b + qn \quad \text{para algunos } k,q \in \Z$. \\
Sumando:
$a' + b' = (a+b) + (k+q)n$. \\

Por lo tanto: \\
$a' + b' \equiv a+b \pmod{n}$. \\
Esto implica: \\
$[a'+b'] = [a+b]$. \\
Producto: \\
Supongamos que tomamos otros representantes $a' \in [a]_n$ y $b' \in [b]_n$. \\  
Eso significa: \\
$a' \equiv a \pmod{n}, \quad b' \equiv b \pmod{n}$. \\
Entonces:
$a' = a + kn, \quad b' = b + qn \quad \text{para algunos } k,q \in \Z$. \\
Efectuando el producto:
$a'b' = (a+kn)(b+q n) = ab + (aq + bk + knq)n$. \\
Por lo tanto: \\
$a' \cdot b' \equiv a \cdot b \pmod{n}$. \\
Esto implica: \\
$[a'\cdot b'] = [a \cdot b]$. \\
Lo que muestra que ambas operaciones están bien definidas.\\

Una vez instituidas las operaciones a través de sus definiciones y mostrado que están bien definidas, ambas expresiones denotan al mismo conjunto, es solo una notación diferente para el mismo objeto. Por lo tanto podemos tratarla como una identidad algebraica: \\
$[a]_n + [b]_n = [a + b]_n$ \\
$[a]_n \cdot [b]_n = [a \cdot b]_n$ \\
El símbolo “=” ya no está introduciendo la definición, sino afirmando que el resultado de aplicar la operación coincide con la clase de equivalencia de la suma (o producto) de representantes. \\
La definición garantiza que la igualdad siempre se cumple. 

\section{Los anillos}
$\Z/n\Z$ \textbf{es un anillo} (el anillo de los enteros módulo $n$): \\
Las operaciones están bien definidas y cumplen los requisitos de anillo.\\
Neutro aditivo: $[0]_n$, pues $[a]_n + [0]_n = [a+0]_n = [a]_n$. \\
Inverso aditivo: $[-a]_n$, pues $[a]_n + [-a]_n = [a-a]_n = [0]_n$. \\
Asociatividad y conmutatividad de la suma: se heredan de los enteros, pues $(a+b)+c = a+(b+c)$ y $a+b = b+a$ en $\Z$, y por lo tanto en las clases. \\
Neutro multiplicativo: $[1]_n$, pues $[a]_n \cdot [1]_n = [a \cdot 1]_n = [a]_n$. \\
Asociatividad del producto: se hereda de $\Z$. \\
Distributividad: se hereda de $\Z$, pues $[a]_n \cdot ([b]_n + [c]_n) = [a]_n \cdot ([b+c]_n) = [a \cdot (b+c)]_n = [(a \cdot b) + (a \cdot c)]_n = [(a \cdot b)]_n + (a \cdot c)]_n =  ([a]_n \cdot [b]_n) + ([a]_n \cdot [c]_n)$. \\
De igual modo $\Z/p\Z$ y $\Z/q\Z$ son anillos. \\

$\Z/p\Z \times \Z/q\Z$ \textbf{es un anillo}: ($\times$ denota el producto directo o producto cartesiano). \\
Es decir, los elemento de $\Z/p\Z \times \Z/q\Z$ son pares ordenados ($[a]_p$,$[b]_q$), el primer componente $[a]_p \in \Z/p\Z$, el segundo componente $[b]_q \in \Z/q\Z$ y la suma y producto son componente a componente (o coordenada a coordenada).\\
Los requisitos (neutros, asociatividad, etc.) se verifican coordenada a coordenada, heredándolos de $\Z/p\Z$ y $\Z/q\Z$. \\

En general, dada la factorización prima $n = p_1^{a_1} p_2^{a_2} \cdots p_k^{a_k}$ \\
$\Z/p_1^{a_1}\Z \times \dots \times \Z/p_k^{a_k}\Z$ es un anillo.

\section{Homomorfismo e isomorfismo de anillos}
\textbf{Homomorfismo}: \\
Dados dos anillos $R$ y $S$, una función $\phi: R \to S$ es un homomorfismo de anillos si respeta ambas operaciones y los neutros multiplicativos: \\
$\forall a, b \in R \quad \phi (a + b) = \phi(a) + \phi(b)$. \\
$\forall a, b \in R \quad \phi (a \cdot b) = \phi (a) \cdot \phi (b)$. \\
$\phi (1_R) = 1_S$. \\
La estructura algebraica se transporta intacta de R a S.\\

\textbf{Isomomorfismo}:
Un homomorfismo de anillos $\phi$ es un isomorfismo si además es biyectivo: \\
Inyectivo: si $\phi (a) = \phi (b)$ entonces $a = b$. No colapsa elementos distintos. \\
Sobreyectivo: $\forall c \in S \quad \exists a \in R$ tal que $\phi (a) = c$. Cubre todo S. \\
Un isomorfismo entre $R$ y $S$ se denota por $R \cong S$ y se dice que $R$ y $S$ son isomorfos.


\section{Grupo multiplicativo y grupo cíclico}
$(\Z/n\Z)^{*}$ \textbf{es un grupo} (el grupo \textbf{multiplicativo} de los enteros módulo $n$): \\
$(\Z/n\Z)^{*} = \{ [a]_n \in \Z/n\Z \;/\; mcd(a,n) = 1\}$.  \\
Es decir, por el Enunciado~\ref{enun:existInvModular}, se toman sólo los elementos invertibles de $\Z/n\Z$. \\
Es claro que $(\Z/n\Z)^{*}$ ya no es un anillo pues no es cerrado bajo la suma, pero el producto es asociativo -se hereda de los enteros-, tiene neutro (el $[1]_n$) y todo elemento tiene inverso -por la forma de construirlo-, suficiente para denominarlo grupo. \\
Básicamente se toman los representantes canónicos que están en el $CRR(n)$, luego el orden del grupo es $\phi (n)$.  \\
Dado que se denomina unidad a un elemento que tiene inverso multiplicativo, otra forma de describirlo es que $(\Z/n\Z)^{*}$ es el grupo de unidades del anillo $(\Z/n\Z)$. \\

\begin{enunciado}
$(\Z/n\Z)^{*}$ es cerrado bajo el producto.
\end{enunciado}

 \begin{proof}
 Sean $[a]_n,[b]_n \in (\Z/n\Z)^{*}$. \\
 Por definición,  $[a]_n \cdot [b]_n = [a \cdot b]_n$. \\
 Por definición, $mcd(a,n)=1$ y $mcd(b,n)=1$. \\
 Por el Enunciado~\ref{enun:maspropMcd}.iii, $mcd(a \cdot b,n)=1.$
 \end{proof}

Recordemos que podemos identificar por comodidad la clase a través de su representante canónico.
Cuando $p$ es primo (impar, para nuestros propósitos), el grupo multiplicativo $(\Z/n\Z)^{*}$ agrega una característica: \\
Existe al menos un elemento especial (llamado generador o raíz primitiva) que, al elevarlo a potencias sucesivas, produce todos los números del conjunto antes de regresar al $1$. \\
Cuando ello sucede, recibe el nombre de \textbf{grupo cíclico}. \\
Aunque no lo haremos, en general se puede probar que, para un grupo cíclico de orden $n$, hay $\phi(\phi(n))$ generadores.

\begin{ejemplo}
Sea $p=7$.
Con la operación producto, el grupo cíclico: \\
$(\Z/7\Z)^{*} = \{ [1]_7,[2]_7,[3]_7,[4]_7,[5]_7,[6]_7 \} = \{ 1,2,3,4,5,6 \}$
tiene dos generadores, 3 y 5. \\
$3^1 = 3$ \\
$3^2 = 9 \equiv 2 \pmod 7$ \\
$3^3 = 6 \equiv 6 \pmod 7$ \\
$3^4 = 18 \equiv 4 \pmod 7$ \\
$3^5 = 12 \equiv 5 \pmod 7$ \\
$3^6 = 15 \equiv 1 \pmod 7$ 
\end{ejemplo}

Para cualquier entero $a$ (que no sea múltiplo de $p$, pues en ese caso es congruente con $0$), se llama orden de $a$ al menor entero positivo $d$ tal que: \\
$a^d \equiv 1 \pmod p$
En el ejemplo, el orden del generador $3$ es $6$. \\
El orden de $37$ es $3$, pues $37^1 \mod 7 = 2, 37^2 \mod 7 = 4, 37^3 \mod 7 = 1$. \\
Si recordamos el pequeño teorema de Fermat, el orden del generador $g$ (que al ser un elemento del grupo es coprimo con $p$) siempre es $p-1$: $g^{p-1} \equiv 1 \pmod p$. \\

Los enunciados de Fermat y Miller-Rabin, entre otros, se basan en la coprimalidad, es decir, en un grupo multiplicativo. \\
Cuando se apela a generadores, estamos trabajando en un grupo multiplicativo cíclico, por ejemplo cuando se descomponen las condiciones de un mentiroso en un sistema de ecuaciones. \\
El teorema del resto chino, convertir ecuaciones multiplicativas en lineales y un largo etcétera ponderan la importancia de un grupo multiplicativo (cíclico, si es el caso).

\section{Homomorfismo e isomorfismo de grupos. \\ 
El isomorfismo de exponenciación discreta}
Sean $(G,+)$ y $(H,\cdot)$ dos grupos. \\
Una función $\varphi:G\to H$ es un isomorfismo de grupos si: \\
1. $\varphi$ es biyectiva. \\
2. $\varphi(a+b)=\varphi(a)\cdot\varphi(b)$ para todos $a,b\in G$. \\

No es difícil ver que $(\Z/(p-1)\Z,+)$ es un grupo.\\
$\Z/(p-1)\Z = \{0,...,p-2\}$. También es equivalente usar $\{1,...,p-1\}$. \\
La operación es la suma módulo $p-1$. \\
La asociatividad y el neutro se heredan de los enteros. \\
$-a$ representa el inverso conceptual de la clase cuyo representante es $a$, aunque una forma de dar un representante positivo es decir que el inverso de $a$ es $p-1-a$.\\

Ya vimos que $(\Z/p\Z)^{*}$ es un grupo. \\
Podemos anotarlo como $(\Z/p\Z,*)$ enfatizando la operación producto.

\begin{enunciado} [El isomorfismo exponencial]\label{enun:isoexpo}\leavevmode\par
Dado $p$ un primo impar: $(\Z/(p-1)\Z,+) \cong (\Z/p\Z)^{*}$ \\
Es decir, $(\Z/(p-1)\Z,+)$ es isomorfo a $(\Z/p\Z)^{*}$.
\end{enunciado}

\begin{proof} \textbf{Parte 1}. \\
Definimos la función $\varphi : (\Z/(p-1)\Z,+) \to (\Z/p\Z)^{*}$: \\
$\varphi ([t]_{p-1}) = [g^t]_p$. \\
Ya advertimos que todo puede hacerse con clases aunque es absolutamente estándar trabajar sólo con los representantes canónicos (los módulos están implícitos): \\
$\varphi (t) = g^t$. 
\end{proof}

Se hace un paréntesis en la prueba para dejar sentados algunos apuntes.\\
La función toma un exponente (al que se eleva el generador de $(\Z/p\Z)^{*}$) y produce un valor en dicho grupo multiplicativo cíclico. \\
Como $(\Z/p\Z)^{*}$ tiene $p-1$ elementos, se requieren $p-1$ exponentes t distintos tales que $g^t$ obtenga tales elementos. \\
Es por ello que se trabaja con el grupo $(\Z/(p-1)\Z,+)$, donde están los exponentes. \\
Con menos exponentes no cubriríamos el orden de $(\Z/p\Z)^{*}$, con más algunos de ellos serían redundantes. \\
En general, $g^t = g^{(t \mod (p-1))}$. \\

Salvo los múltiplos de $p$, todos los demás enteros $a$ están en alguna clase de $(\Z/p\Z)^{*}$. Que el generador $g$ genere $a$ quiere decir que genera al representante de la clase a la que $a$ pertenece.

\begin{ejemplo} 
Con $p=7$ el entero $a=5555$ está en la misma clase que el $4$. Luego, con el generador $g=3$, existe un $t$ en $(\Z/(p-1)\Z,+)$ tal que $g^t \equiv 4 \pmod 7$, en efecto, $3^4 \equiv 4 \pmod 7$. \\
Para cualquier $p$ primo impar: $g^0 \equiv g^{p-1} \equiv 1 \pmod p$. \\
Con $p=11$ el grupo multiplicativo tiene estos elementos (representantes): $\{1, \dots, 10\}$. \\
Para generar los enteros $a=10$ y $a=1$ con el generador $g=8$, elegimos un exponente $t$ en $\{0,...,9\}$ o, equivalentemente, en $\{1,...,10\}$: \\
$8^{5} \equiv 10 \pmod {11}$, $t=5$. \\
$8^{10} \equiv 1 \pmod {11}$, $t=10$, también sirve $t=0$.
\end{ejemplo}

Retomamos la prueba.

\begin{proof} \textbf{Parte 2}. \\
$\varphi$ es un homomorfismo: \\
$\varphi(t+t') = g^{t+t'} = g^t\cdot g^{t'} = \varphi(t)\cdot\varphi(t')$. \\
$\varphi$ \textbf{está bien definida}: \\
Sea $t\equiv t'\pmod {p-1}$. \\
Entonces por el Enunciado~\ref{enun:congrViaDivisib}, $(p-1) \mid (t-t')$, es decir, $t'=t+(p-1)q$ para algún entero $q$. \\
Como $g^{p-1} \equiv 1 \pmod p$, se tiene $g^{t'}=g^t \cdot (g^{p-1})^q \equiv g^t\pmod p$. \\
Es decir, $\varphi(t)$ no depende del representante elegido en $\Z/n\Z$. \\
$\varphi$ \textbf{es inyectiva}: \\
Si $\varphi(t)=\varphi(t')$, es decir, $g^t\equiv g^{t'}\pmod p$, entonces $g^{t-t'}\equiv1 \pmod p$. \\
Por el algoritmo de la división: $t-t' = q(p-1)+r$, con $0 \leq r < (p-1)$. Luego: \\
$(g^{p-1})^q \cdot g^r \equiv 1 \pmod p$. \\
Como el orden del generador $g$ es $p-1$ (o por el pequeño teorema de Fermat), resulta: \\
$g^{p-1} \equiv 1 \pmod p$. \\
Por el Enunciado~\ref{enun:propCongru}.viii resulta: $g^r \equiv 1 \pmod p$. \\
Si $r$ fuera positivo, significaría que el orden de g no es $p-1$, por tanto $r=0$. \\
Entonces, $t-t' = q(p-1)$, es decir, $(p-1) \mid (t-t')$. Por lo tanto: \\ 
$t \equiv t' \pmod {p-1}$. \\
Así $\varphi$ no colapsa elementos distintos de $\Z/(p-1)\Z$. \\
$\varphi$ \textbf{es sobreyectiva}: \\
$\forall c \in (Z/pZ)^{*}$ existe un exponente $t$ tal $\varphi(t) = g^t = c$. \\
Es precisamente la definición de generador $g$.
\end{proof}

Lo que hace $\varphi$ es convertir una suma (del mundo aditivo de exponentes) en un producto (del mundo multiplicativo de residuos).\\

Dado que $\varphi$ está bien definida: $t\equiv t' \pmod {p-1} \implies g^t=g^{t'}$. \\
Dada la inyectividad de $\varphi$: $g^t=g^{t'} \implies t\equiv t' \pmod {p-1}$. \\
Lo que nos permite escribir el siguiente enunciado:

\begin{enunciado} [El isomorfismo exponencial de grupos]\label{enun:isoexpo2}\leavevmode\par
$g^t=g^{t'} \iff t\equiv t' \pmod {p-1}$
\end{enunciado}

\section{Ecuación lineal en congruencia}
Sea $x$ la incógnita de la ecuación lineal en congruencia: \\
$ax \equiv b \pmod {n}$. \\
Sea $d = mcd(a,n)$. 

\begin{enunciado} [Número de soluciones de una ecuación lineal en congruencia]\label{enun:numsolineal}\leavevmode\par
$ax \equiv b \pmod {n}$ tiene solución si y solo si $d \mid b$. \\
En ese caso tiene exactamente $d$ soluciones distintas módulo $n$.
\end{enunciado}

\begin{proof} \, \\
$\Rightarrow$) \\
Supongamos que $x_0$ es solución de $ax\equiv b \pmod n$, es decir: \\
$ax_0 \equiv b \pmod n$. \\
Por el Enunciado~\ref{enun:congrViaDivisib}: $n \mid (ax_0-b)$. \\
Esto significa que existe un entero $q$ tal que
$ax_0 - b = qn $, es decir, $b = ax_0 - qn$. \\
Como $d = mcd(a,n)$, es obvio que $d \mid a$ y que $d \mid n$. \\
Por el Enunciado~\ref{enun:propDivisibilidad}: $d \mid ax_0-qn$. Luego: \\
$d\mid b$. \\
$\Leftarrow$) \\
Supongamos $d\mid b$. \\
Como $d = mcd(a,n)$, por la identidad de Bezout, Enunciado~\ref{enun:propMcd}.i, existen enteros $u,v$ tales que: \\
$d = au + nv$. \\
Como $d\mid b$, entonces $b = dc$ para algún entero $c$. \\
Multiplicando la identidad de Bezout por $c$: \\
$b = dc = auc + nvc$. Es decir: \\
$a(uc) - b = n(-vc)$. Es decir: \\
$ n \mid a(uc) - b$. Es decir, por el Enunciado~\ref{enun:congrViaDivisib}: \\
$a(uc)\equiv b\pmod n$.\\
Por lo tanto, $x=uc$ es una solución explícita. \\
Se muestra a continuación la forma de las soluciones. \\
Sea $x_0$ una solución particular (la que se acaba de construir, o cualquier otra). \\
Cualquier otra solución tiene la forma $x_0 + n'j$ para algún entero $j$, donde $n'=n/d$.\\
En efecto, sea $x'$ otra solución, así que $ax'\equiv b\equiv ax_0\pmod n$, es decir, por el Enunciado~\ref{enun:propCongru}.i y luego por el Enunciado~\ref{enun:congrEquiv}: \\
$a(x'-x_0)\equiv 0\pmod n \iff n\mid a(x'-x_0)$. \\
Como $d = mcd(a,n)$, es obvio que $d \mid a$ y que $d \mid n$. Es decir:\\
$a=da'$, $n=dn'$ (aquí resulta que $n'=n/d$). \\
Es claro que $mcd(a',n')=1$, pues si $mcd(a',n')=h$, con $h>1$, como $h \mid a'$ y $h \mid n'$, resulta $a'=hq_1$ y $n'=hq_2$. Sustituyendo, resulta: $a=(dh)q_1$, $n=(dh)q_2$. \\
Como $h > 1$, entonces $dh > d$. Es decir, hay un divisor común de $a$ y $n$ más grande que $d$. Imposible.  \\
Entonces la condición $n \mid a(x'-x_0)$ se escribe $dn' \mid da'(x'-x_0) \iff n' \mid a'(x'-x_0)$. \\
Como $mcd(a',n')=1$, por el lema de Euclides, Enunciado~\ref{enun:propMcd}
.viii: $n' \mid (x'-x_0)$, 
es decir, $x' = x_0 + n' j$ para algún entero $j$. \\
Recíprocamente, cualquier $x'$ de esta forma es solución: \\
$a(x_0+n'j) = ax_0 + an'j = ax_0 + (da')n' j = ax_0 + a'\,(dn')\,j = ax_0 + a' n j\equiv ax_0\equiv b\pmod n$. \\
Ahora sí contemos las soluciones distintas módulo $n$. \\
Las soluciones son $x_0, x_0+n', x_0+2n'$, etc. Es decir, $x_0 + n'j$ para $j \in \Z$. \\
Cuándo dos de ellas, $x_0+n'j_1$ y $x_0+n'j_2$, son congruentes módulo $n$: \\
$x_0+n'j_1 \equiv x_0+n'j_2 \pmod n \iff n'j_1 \equiv n'j_2 \pmod n \iff n \mid n'(j_1-j_2)$ \\
$\iff dn'\mid n'(j_1-j_2) \iff d\mid (j_1-j_2) \iff j_1 \equiv j_2 \pmod d$. \\
Por lo tanto, reemplazando los valores $j=0,1,2,\dots,d-1$ dan soluciones de la ecuación módulo $n$ todas distintas, pues tomando estos elementos diferentes por pares, ninguno es congruente módulo $d$, pues representan las $d-1$ clases de residuos módulo $d$. \\
Cualquier otro $j$ es congruente (módulo $d$) a alguno de estos, dando una solución repetida. \\
Por lo tanto hay exactamente $x_0,\; x_0+n',\; x_0+2n',\;\dots,\; x_0+(d-1)n'$ soluciones distintas módulo $n$, es decir, $d$ soluciones.
\end{proof}

\section{Anillo de polinomios, divisibilidad y teorema del factor}
\label{sec:polinomios}
Sea $R$ un anillo conmutativo con unidad. 

\begin{definicion}[Anillo de polinomios]\label{def:anilloPolinomios}
El anillo de polinomios denotado por $R[x]$ es el conjunto de todas las expresiones formales $a_0 + a_1x + a_2x^2 + \cdots + a_nx^n$ \\
donde: \\
Los coeficientes $a_i$ pertenecen al anillo $R$. \\
El símbolo $x$ es una indeterminada (no un número fijo). \\
El grado del polinomio $n$ es finito (cada polinomio tiene un número finito de términos). \\
\textbf{Operaciones} \\
En $R[x]$ se definen dos operaciones naturales. \\
Dados \\
$f(x) = a_0 + a_1x + a_2x^2 + \cdots + a_nx^n \quad \operatorname{gr}(f)=n$ \\
$g(x) = b_0 + b_1x + b_2x^2 + \cdots + b_mx^m \quad \operatorname{gr}(g)=m$ \\
Suma de polinomios (se usa la suma de $R$): \\
$f(x) + g(x) = (a_0+b_0) + (a_1+b_1)x + (a_2+b_2)x^2 + \cdots + (a_k+b_k)x^k \qquad$ $k = \max(n,m)$ \\
Si algún índice excede el grado de uno de los polinomios, se toma el coeficiente faltante como 0. \\
Producto de polinomios (se usa el producto de $R$): \\
$f(x) \cdot g(x) = c_0 + c_1x + c_2x^2 + \cdots + c_{n+m}x^{n+m} \qquad$ $c_k = \sum_{i+j=k} a_i b_j$ \\
\end{definicion}

Nótese que la suma es término a término. \\
Nótese que el producto es distributivo. \\
Si $R$ es un cuerpo, entonces $R[x]$ no es un cuerpo porque no todos los polinomios tienen inverso. \\
$R[x]$ es un dominio de integridad,

\begin{definicion}[Dominio de integridad]\label{def:domDeIntegr}
Un dominio de integridad es un anillo conmutativo con unidad sin divisores de cero, es decir, tiene suma y producto, el producto es conmutativo y existe un elemento neutro multiplicativo $1$. \\
La condición de no tener divisores de cero es esta: \\  
Si $a \cdot b = 0$ entonces $a = 0$ o $b = 0$, para cualesquier $a,b$ en el anillo.
\end{definicion}

En el producto de polinomios, con $a_n \neq 0$ y $b_m \neq 0$, el término de mayor grado de $f(x) \cdot g(x)$ es $a_n b_m x^{n+m}$, como $R$ es un dominio de integridad, $a_n \cdot b_m \neq 0$, por lo tanto el producto no se anula, no hay divisores de cero, \\

En un cuerpo, dado que todo elemento distinto de cero tiene inverso multiplicativo, se puede probar directamente que no hay divisores de cero. \\
En efecto, supongamos que en un cuerpo hubiera divisores de cero: $a \neq 0, b \neq 0$, pero $a \cdot b = 0$. \\
Como $a \neq 0$, existe $a^{-1}$. Multiplicando por $a^{-1}$: \\  
$a^{-1} \cdot (a \cdot b) = a^{-1} \cdot 0$, entonces $(a^{-1} \cdot a) \cdot b = 0$, luego $1 \cdot b = 0$, es decir, $b = 0$. \\
Lo que contradice que $b \neq 0$, invalida el supuesto y prueba la afirmación. \\

$\Z/8\Z$ (un anillo conmutativo con unidad), por ejemplo, tiene divisores de cero: $2 \cdot 4$. \\
En efecto, $2 \cdot 4 \bmod 8 = 0$, aún cuando $2 \neq 0$ y $4 \neq 0$. \\

Existe un algoritmo de la división para polinomios.

\begin{enunciado}[Algoritmo de la División de polinomios]\label{enun:algDivisionPol}
Sea $K$ un cuerpo y $p(x), d(x) \in K[x]$ con $d(x) \neq 0$. Entonces existen polinomios $q(x), r(x) \in K[x]$ tales que: \\
$p(x) = d(x)q(x) + r(x) \qquad$ $\operatorname{gr}(r) < \operatorname{gr}(d)$ \\
$q(x)$ y $r(x)$ son únicos y se denominan el cociente y el residuo.
\end{enunciado}

\begin{proof}
Existencia \\
Por inducción sobre $n = \operatorname{gr}(p)$ \\
Sea $m = \operatorname{gr}(d)$ y $b$ el coeficiente principal de $d(x)$. \\
Caso base: Sea $\operatorname{gr}(p) < m$. \\
Tomamos $q(x) = 0$ y $r(x) = p(x)$. Se cumple trivialmente. \\
Hipótesis inductiva: El enunciado es cierto para $\operatorname{gr}(p) < n$. \\
Paso inductivo: Supongamos $\operatorname{gr}(p) = n \geq m$. \\
Sea $a$ el coeficiente principal de $p(x)$. Definimos: $p_1(x) = p(x) - \frac{a}{b} x^{n-m} d(x)$. \\
Como restamos el término principal, $\operatorname{gr}(p_1) < n$. Por hipótesis inductiva, existen $q_1(x), r(x)$ tales que: $p_1(x) = d(x)q_1(x) + r(x), \quad \operatorname{gr}(r) < m$. \\
Sustituyendo: \\
$p(x) - \frac{a}{b} x^{n-m} d(x) = d(x)q_1(x) + r(x)$. \\
Luego: \\
$p(x) = d(x) \left(q_1(x) + \frac{a}{b} x^{n-m}\right) + r(x)$. \\
Es decir: \\
$p(x) = d(x) \cdot q(x) + r(x) \quad$ donde $q(x) = \left(q_1(x) + \frac{a}{b} x^{n-m}\right)$. \\
Como $b$ es el coeficiente principal de $d(x)$ y $d(x) \neq 0$, entonces $b \neq 0$. \\
En el cuerpo $K$ todo elemento no nulo tiene inverso multiplicativo, por lo tanto, existe $b^{-1} \in K$. Luego, $\frac{a}{b} = a \cdot b^{-1} \in K$. \\
Unicidad \\
Supongamos dos representaciones: \\
$p(x) = d(x)q_1(x) + r_1(x) = d(x)q_2(x) + r_2(x)$ \\
Restando: $d(x)(q_1(x) - q_2(x)) = r_2(x) - r_1(x) \quad$  (lado-izquierdo LI = LD lado-derecho)\\
Supongamos que $q_1 \neq q_2$. \\
Entonces $\operatorname{gr}(\text{LI}) \geq \operatorname{gr}(d) = m$. \\
Pero $\operatorname{gr}(\text{LD}) \leq \max(\operatorname{gr}(r_1), \operatorname{gr}(r_2)) < m$. \\
Esta contradicción invalida el supuesto, por lo tanto, $q_1 = q_2$, y en consecuencia $r_1 = r_2$.
\end{proof}

\begin{enunciado}[Teorema del residuo]\label{enun:teoResid}
Sea $p(x)  \in K[x]$. \\
Sea $c \in K$. \\
Si se divide $p(x)$ por el polinomio $(x - c)$ entonces el resto $r(x)$ de dicha división es el elemento constante $r' \in K$ tal que: $r' = p(c)$
\end{enunciado}

\begin{proof}
Por el Algoritmo de la División de polinomios, existen únicos polinomios $q(x)$ y $r(x)$ tales que: $p(x) = (x - c) \cdot q(x) + r(x) \quad$ con 
$\operatorname{gr}(r) < \operatorname{gr}(x - c) \qquad$  \hfill (*) \\
Es evidente que el grado de $(x - c)$ es $1$. Por lo tanto $r(x)$ es el polinomio constante (de grado 0), es decir, una constante, llamémosla $r'$: $p(x) = (x - c) \cdot q(x) + r'$. \\
Evaluando (*) en $x = c$ resulta: $p(c) = (c - c) \cdot q(c) + r'$, es decir, $p(c) = r'$.
\end{proof}

\begin{enunciado}[Teorema del factor]\label{enun:teoFactor}
Sea $K$ un cuerpo. \\
Sea $p(x) \in K[x]$. \\
Para cualquier elemento $c \in K$, se cumple que: $\,(x - c) \mid p(x) \, \iff \, p(c) = 0$, \\
Es decir, $(x - c)$ es un factor de $p(x)$ si y solo si $c$ es una raíz de $p(x)$.
\end{enunciado}

\begin{proof}
Dados los polinomios $p(x), d(x) \in K[x]$, $d(x) \neq 0$, por el Algoritmo de la División para polinomios existen únicos polinomios $q(x), r(x)$ tales que: \\
$p(x) = d(x) \cdot q(x) + r(x), \quad$ con $\operatorname{gr}(r) < \operatorname{gr}(d)$. \\
Aplicado a $d(x) = (x - c)$, por el teorema del residuo: \\
$p(x) = (x - c) \cdot q(x) + r' \quad$ con $r'= p(c)$. \\
$\Leftarrow$) \\
Si $p(c) = 0$ entonces $(x - c) \mid p(x)$. \\
Por hipótesis: $p(c) = 0$. Luego, $r'= 0$. Es decir: \\
$p(x) = (x - c) \cdot q(x)$. Lo que implica que: \\
$(x - c) \mid p(x)$. \\
$\Rightarrow$) \\
Si $(x - c) \mid p(x)$ entonces $p(c) = 0$. \\
Por hipótesis: $(x - c) \mid p(x)$. Es decir: $p(x) = (x - c) \cdot q(x)$ \\
Evaluando en $x = c$: $p(c) = (c - c) \cdot q(c) = 0$. Es decir:\\
$p(c) = 0$.
\end{proof}

Nótese que la demostración es la misma para los polinomios con coeficientes en el cuerpo $Z/pZ$ (obviamente las operaciones, igualdades y evaluaciones se entienden módulo $p$).

\begin{enunciado}[Número máximo de raíces]\label{enun:numRaices}
Sea $K$ un cuerpo. \\
Todo polinomio no nulo $p(x)\in K[x]$ de grado $n$ tiene a lo sumo $n$ raíces distintas en $K$
\end{enunciado}

\begin{proof}
Por inducción sobre $n$. \\
Caso base: $n=0$ \\
Un polinomio no nulo de grado $0$ es una constante: $p(x)=c$,\, con $c\neq0$. \\
Por tanto no tiene raíces, es decir, tiene $0$ raíces: $0 \leq n$. \\
Hipótesis inductiva: el enunciado es cierto para todo polinomio no nulo de grado $n$. \\
Paso inductivo: Sea $p(x)$ un polinomio no nulo de grado $n+1$. \\
Si $p(x)$ no tiene raíces, entonces tiene $0$ raíces, y evidentemente $0\leq n+1$. \\
Si $p(x)$ tiene alguna raíz $r' \in K$, quiere decir que $p(r')=0$. Por el teorema del factor, $x-r' \mid p(x)$. \\
Por tanto, existe $g(x) \in K[x]$ tal que $p(x)=(x-r')g(x)$. \\
Como $\operatorname{gr}(p)=n+1$, se deduce que $\operatorname{gr}(g)=n$, \\
Todas las demás raíces de $p(x)$, distintas de $r'$, son raíces de $g$. \\
En efecto, si $s \neq r'$ es otra raíz de $p(x)$, entonces $0=p(s)=(s-r)g(s)$. \\
Como $K$ es un cuerpo, no tiene divisores de cero. Como $s \neq r'$, $s-r' \neq 0$. \\
Por tanto necesariamente $g(s)=0$. \\
Así que: $r'$ es una raíz de $p(x)$; todas las demás raíces de $p(x)$ son raíces de $g(x)$; $g(x)$ tiene grado $n$. \\
Por la hipótesis inductiva, $g(x$) tiene a lo sumo $n$ raíces. \\
Por consiguiente, $p(x)$ tiene como máximo $1+n=n+1$ raíces. \\
Los tres casos demuestran que todo polinomio no nulo de grado $n$ sobre un cuerpo tiene a lo sumo $n$ raíces.
\end{proof}

Considerando  el polinomio $p(x)=x^2-1$ en el cuerpo $\Z/p\Z$, $p(x)$ tiene a lo sumo $2$ raíces. \\
Sabemos que $1$ es raíz porque $p(1)=0$. Por el teorema del factor, $x-1 \mid x^2-1$. \\
De hecho, $x^2-1=(x-1)(x+1)$. \\
Ya tenemos una raíz $x=1$. La otra es $x=-1$. \\
Es decir, $p(x)=x^2-1$ tiene exactamente $2$ raíces.
